% SEGMENT OLP-0024-S01
% Part: sets-functions-relations
% Chapter: functions
% Section: inverses

\documentclass[../../../include/open-logic-section]{subfiles}

\begin{document}

\olfileid{sfr}{fun}{inv}
\olsection{சார்புகளின் நேர்மாறுகள்}

\begin{explain}
சார்புகளை இணைப்புகளாகக் கருதுகிறோம். எனவே ஒரு சார்பு ஏற்படுத்தும்
இணைப்பை “திருப்ப முடியுமா” என்ற இயல்பான கேள்வி எழுகிறது.
எடுத்துக்காட்டாக, $f(x) = x + 1$ என்ற தொடரிச் சார்பு $f$ செய்வதை
$g(y) = y - 1$ என்ற சார்பு “மீட்டுவிடுகிறது” என்ற பொருளில்,
அந்தத் தொடரிச் சார்பைத் திருப்ப முடியும்.

ஆனால் கவனம் தேவை. $g$ இன் வரையறை $\Int \to \Int$ என்ற
ஒரு சார்பை வரையறுத்தாலும், அது $\Nat
\to \Nat$ என்ற ஒரு \emph{சார்பை} வரையறுக்கவில்லை.
ஏனெனில் $g(0) \notin \Nat$.
ஆகவே எளிய எடுத்துக்காட்டுகளில்கூட ஒரு சார்பைத் திருப்ப முடியுமா
என்பது முற்றிலும் வெளிப்படையானதல்ல; அது சார்பகத்தையும்
துணைச்சார்பகத்தையும் பொறுத்திருக்கலாம்.

ஒரு சார்பின் நேர்மாறு என்ற கருத்து இதனை மேலும் துல்லியமாக்குகிறது.
\end{explain}

% SEGMENT OLP-0024-S02
\begin{defn}
அனைத்து $x \in A$, $y \in B$ ஆகியவற்றுக்கும்
$f(g(y)) = y$, $g(f(x)) = x$ என அமையுமானால்,
$g \colon B \to A$ என்ற சார்பு, $f
\colon A \to B$ என்ற சார்பின் ஒரு \emph{நேர்மாறு} ஆகும்.
\end{defn}

$f$ க்கு ஒரு நேர்மாறு $g$ இருந்தால், $g$ என்பதற்குப் பதிலாக
அடிக்கடி $f^{-1}$ என எழுதுகிறோம்.

% SEGMENT OLP-0024-S03
\begin{explain}
சார்புகளுக்கு எப்போது நேர்மாறுகள் உள்ளன என்பதை இப்போது காண்போம்.
$f\colon A \to B$ இன் நேர்மாறாக அமையக்கூடிய ஒரு நல்ல தேர்வு,
பின்வருமாறு “வரையறுக்கப்படும்” $g\colon B \to A$ ஆகும்:
\[
g(y) = \text{“அந்த” $x$; இங்கு $f(x) = y$.}
\]
ஆனால் “வரையறுக்கப்படும்”, “அந்த” என்பவற்றைச் சுற்றியுள்ள மேற்கோள்
குறிகள், இது இன்னும் ஒரு வரையறை அல்ல என்பதைக் காட்டுகின்றன.
குறைந்தபட்சம், முழுப் பொதுமையுடன் இது எப்போதும் செயல்படாது.
ஏனெனில் இந்த வரையறை ஒரு சார்பைக் குறிப்பிட வேண்டுமானால்,
$f(x) =
y$ என அமையும் ஒரே ஓர் $x$ இருக்க வேண்டும்:
$g$ இன் வெளியீடு தனித்ததாகக் குறிப்பிடப்பட வேண்டும்.
மேலும், ஒவ்வொரு $y \in B$ க்கும் அது குறிப்பிடப்பட வேண்டும்.
$x_1$, $x_2 \in
A$ ஆகியவை $x_1 \neq x_2$ எனவும், ஆனால்
$f(x_1) = f(x_2)$ எனவும் இருந்தால்,
$y = f(x_1) = f(x_2)$ என்பதற்கு $g(y)$ தனித்ததாகக்
குறிப்பிடப்படாது. மேலும், $f(x) = y$ என அமையும்
$x$ எதுவுமே இல்லை எனில், $g(y)$ குறிப்பிடப்படவே இல்லை.
வேறு சொற்களில், $g$ வரையறுக்கப்பட வேண்டுமானால்,
$f$ என்பது !!{injective}, !!{surjective} ஆகிய இரு பண்புகளையும்
கொண்டிருக்க வேண்டும்.

படிப்படியாகப் பார்ப்போம். கேள்வியை இரண்டாகப் பிரிப்போம்.
$f\colon A \to B$ என்ற சார்பு தரப்பட்டிருக்கும்போது,
$g(f(x)) = x$ என அமையும் $g\colon B \to A$ என்ற சார்பு
எப்போது உள்ளது? இத்தகைய $g$, $f$ செய்வதை “மீட்டுவிடுகிறது”;
அது $f$ இன் \emph{இடது நேர்மாறு} எனப்படும்.
இரண்டாவதாக, $f(h(y)) = y$ என அமையும்
$h\colon B \to A$ என்ற சார்பு எப்போது உள்ளது?
இத்தகைய $h$, $f$ இன் \emph{வலது நேர்மாறு} எனப்படும்:
$h$ செய்வதை $f$ “மீட்டுவிடுகிறது”.
\end{explain}

% SEGMENT OLP-0024-S04
\begin{prop}
சார்பகம் வெற்றுக் கணமல்ல எனவும் $f\colon A \to B$ என்பது
!!{injective} பண்புடையது எனவும் இருந்தால்,
அனைத்து $x
\in A$ க்கும் $g(f(x)) = x$ என அமையும்
$f$ இன் \emph{இடது நேர்மாறு} $g\colon B \to A$ ஒன்று உள்ளது.
\end{prop}

% SEGMENT OLP-0024-S05
\begin{proof}
சார்பகம் வெற்றுக் கணமல்ல எனவும் $f\colon A \to B$ என்பது
!!{injective} பண்புடையது எனவும் கொள்வோம்.
ஒரு $y \in B$ ஐக் கருதுவோம்.
$y \in \ran{f}$ எனில், $f(x) = y$ என அமையும்
ஒரு $x \in A$ உள்ளது. $f$ என்பது !!{injective} பண்புடையதால்,
இத்தகைய $x \in A$ ஒன்றே உள்ளது. எனவே
$g(y) = x$ என வரையறுக்கலாம்; அதாவது,
$g(y)$ என்பது $f(x)
= y$ என அமையும் “அந்த” $x \in A$ ஆகும்.
$y \notin \ran{f}$ எனில், அதை ஏதேனும் ஓர் $a \in A$ உடன்
இணைக்கலாம். ஆகவே ஓர் $a \in A$ ஐத் தேர்ந்தெடுத்து,
$g\colon B \to A$ என்பதைப் பின்வருமாறு வரையறுக்கலாம்:
\[
g(y) = \begin{cases}
    x & \text{$f(x) = y$ எனில்}\\
    a & \text{$y \notin \ran{f}$ எனில்.}
\end{cases}
\]
இது அனைத்து $y \in B$ க்கும் வரையறுக்கப்பட்டுள்ளது.
ஏனெனில் இத்தகைய ஒவ்வொரு $y \in \ran{f}$ க்கும்,
$f(x) = y$ என அமையும் ஒரே ஓர் $x \in A$ உள்ளது.
வரையறையின்படி, $y = f(x)$ எனில் $g(y) = x$;
அதாவது $g(f(x)) = x$.
\end{proof}

% SEGMENT OLP-0024-S06
\begin{prob}
$f\colon A \to B$ க்கு ஓர் இடது நேர்மாறு $g$ இருந்தால்,
$f$ என்பது !!{injective} பண்புடையது எனக் காட்டுக.
\end{prob}

% SEGMENT OLP-0024-S07
\begin{prop}
$f\colon A \to B$ என்பது !!{surjective} பண்புடையது எனில்,
அனைத்து $y \in B$ க்கும் $f(h(y)) =
y$ என அமையும் $f$ இன் \emph{வலது நேர்மாறு}
$h\colon B \to A$ ஒன்று உள்ளது.
\end{prop}

% SEGMENT OLP-0024-S08
\begin{proof}
$f\colon A \to B$ என்பது !!{surjective} பண்புடையதாகக் கொள்வோம்.
ஒரு $y \in
B$ ஐக் கருதுவோம். $f$ என்பது !!{surjective} பண்புடையதால்,
$f(x_y)
= y$ என அமையும் ஒரு $x_y \in A$ உள்ளது.
எனவே $h(y) = x_y$ என வரையறுக்கலாம்; அதாவது,
ஒவ்வொரு $y \in B$ க்கும் $f(x) = y$ என அமையும்
ஏதேனும் ஓர் $x \in A$ ஐத் தேர்ந்தெடுக்கிறோம்.
$f$ என்பது !!{surjective} பண்புடையதால், தேர்ந்தெடுப்பதற்குக்
குறைந்தபட்சம் ஒன்று எப்போதும் உள்ளது.\footnote{$f$ என்பது
!!{surjective} பண்புடையதால், ஒவ்வொரு $y \in B$ க்கும்
$\Setabs{x}{f(x) = y}$ என்ற கணம் வெற்றுக் கணமல்ல.
$h$ பற்றிய நமது வரையறை, இந்தக் கணங்கள் ஒவ்வொன்றிலிருந்தும்
ஒரே ஒரு $x$ ஐத் தேர்ந்தெடுக்க வேண்டியுள்ளது.
இது எப்போதும் சாத்தியம் என்பது உண்மையில் வெளிப்படையானதல்ல:
இந்தத் தேர்வுகளைச் செய்ய முடியும் என்பது ஓர் அடிக்கோளாகவே
ஏற்றுக்கொள்ளப்படுகிறது. வேறு சொற்களில், இந்த முன்மொழிவு
தேர்வு அடிக்கோள் எனப்படுவதை முன்கொள்கிறது.
இந்தச் சிக்கலை \oliflabeldef{sth:choice::chap}{\olref[sth][choice][]{chap}
இல் மீண்டும் ஆராய்வோம்}{இங்கு விரிவாக ஆராயாமல் கடந்து செல்கிறோம்}.
ஆயினும், பல குறிப்பிட்ட நிலைகளில், எடுத்துக்காட்டாக
$A = \Nat$ ஆகவோ முடிவுறு கணமாகவோ இருக்கும்போது,
அல்லது $f$ என்பது !!{bijective} பண்புடையதாக இருக்கும்போது,
தேர்வு அடிக்கோள் தேவையில்லை.
(குறிப்பாக $f$ என்பது !!{bijective} பண்புடைய நிலையில்,
ஒவ்வொரு $y \in B$ க்கும் $\Setabs{x}{f(x) = y}$ என்ற கணத்தில்
ஒரே ஓர் !!{element} உள்ளது; ஆகவே தேர்ந்தெடுப்பதற்கு
மாற்று உறுப்புகள் ஏதும் இல்லை.)}
வரையறையின்படி, $x = h(y)$ எனில் $f(x) = y$;
அதாவது எந்த $y \in B$ க்கும் $f(h(y)) = y$.
\end{proof}

% SEGMENT OLP-0024-S09
\begin{prob}
$f\colon A \to B$ க்கு ஒரு வலது நேர்மாறு $h$ இருந்தால்,
$f$ என்பது !!{surjective} பண்புடையது எனக் காட்டுக.
\end{prob}

% SEGMENT OLP-0024-S10
\begin{explain}
முந்தைய நிறுவலில் உள்ள கருத்துகளை இணைப்பதால்,
ஒவ்வொரு !!{bijection} க்கும் ஒரு நேர்மாறு உள்ளது என இப்போது பெறுகிறோம்.
அதாவது, $f$ இன் இடது நேர்மாறாகவும் வலது நேர்மாறாகவும்
அமையும் ஒரே சார்பு ஒன்று உள்ளது.
\end{explain}

% SEGMENT OLP-0024-S11
\begin{prop}\ollabel{prop:bijection-inverse}
$f\colon A \to B$ என்பது !!{bijective} பண்புடையது எனில்,
அனைத்து $x \in A$ க்கும் $f^{-1}(f(x)) = x$ எனவும்,
அனைத்து $y \in B$ க்கும் $f(f^{-1}(y)) = y$ எனவும் அமையும்
$f^{-1}\colon B \to A$ என்ற சார்பு ஒன்று உள்ளது.
\end{prop}

% SEGMENT OLP-0024-S12
\begin{proof}
பயிற்சி.
\end{proof}

% SEGMENT OLP-0024-S13
\begin{prob}
\olref[sfr][fun][inv]{prop:bijection-inverse} ஐ நிறுவுக.
நீங்கள் $f^{-1}$ ஐ வரையறுத்து, அது ஒரு சார்பு எனக் காட்டி,
அது $f$ இன் நேர்மாறு எனவும் காட்ட வேண்டும்;
அதாவது, அனைத்து $x \in A$, $y \in B$ ஆகியவற்றுக்கும்
$f^{-1}(f(x)) = x$, $f(f^{-1}(y)) = y$ எனக் காட்ட வேண்டும்.
\end{prob}

% SEGMENT OLP-0024-S14
\begin{explain}
நேர்மாறுகளைப் பெறுவதற்குச் சற்றே பொதுவான மற்றொரு வழியும் உள்ளது.
அனைத்து $x \in A$ க்கும் $f'(x) = f(x)$ என அமைப்பதால்,
ஒவ்வொரு சார்பு $f$ உம் $f'
\colon A \to \ran{f}$ என்ற !!a{surjection} ஐ உருவாக்குகிறது
என்பதை \olref[kin]{sec} இல் கண்டோம்.
தெளிவாக, $f$ என்பது !!{injective} பண்புடையது எனில்,
$f'$ என்பது !!{bijective} பண்புடையதாகும்;
எனவே \olref{prop:bijection-inverse} இன்படி அதற்குத் தனித்த
நேர்மாறு உள்ளது. குறியீட்டை மிகச் சிறிய அளவில் தளர்த்திப் பயன்படுத்தி,
சில சமயங்களில் $f'$ இன் நேர்மாறை வெறுமனே
“$f$ இன் நேர்மாறு” என அழைக்கிறோம்.
\end{explain}

% SEGMENT OLP-0024-S15
\begin{prop}\ollabel{prop:left-right}%
$f\colon A \to B$ க்கு ஓர் இடது நேர்மாறு $g$ உம்
ஒரு வலது நேர்மாறு $h$ உம் இருந்தால், $h = g$ எனக் காட்டுக.
\end{prop}

% SEGMENT OLP-0024-S16
\begin{proof}
பயிற்சி.
\end{proof}

% SEGMENT OLP-0024-S17
\begin{prob}
\olref[sfr][fun][inv]{prop:left-right} ஐ நிறுவுக.
\end{prob}

% SEGMENT OLP-0024-S18
\begin{prop}\ollabel{prop:inverse-unique}
ஒவ்வொரு சார்பு $f$ க்கும் அதிகபட்சம் ஒரே ஒரு நேர்மாறு உள்ளது.
\end{prop}

% SEGMENT OLP-0024-S19
\begin{proof}
$g$, $h$ ஆகிய இரண்டும் $f$ இன் நேர்மாறுகள் எனக் கொள்வோம்.
அப்போது குறிப்பாக, $g$ என்பது $f$ இன் இடது நேர்மாறு;
$h$ என்பது வலது நேர்மாறு.
\olref{prop:left-right} இன்படி, $g = h$.
\end{proof}

\end{document}
